% Generated draft: requires mathematical and source review.
\documentclass[11pt]{article}
\usepackage[margin=25mm]{geometry}
\usepackage{fontspec}
\setmainfont{Latin Modern Roman}
\usepackage{amsmath,amssymb,mathtools,unicode-math}
\setmathfont{Latin Modern Math}
\usepackage{enumitem,needspace,xurl,hyperref}
\setlength{\parindent}{0pt}
\setlength{\parskip}{5pt}
\setlength{\emergencystretch}{3em}
\newcommand{\N}{\mathbb{N}}
\newcommand{\Z}{\mathbb{Z}}
\newcommand{\Q}{\mathbb{Q}}
\newcommand{\R}{\mathbb{R}}
\newcommand{\C}{\mathbb{C}}
\newcommand{\F}{\mathbb{F}}
\newcommand{\A}{\mathbb{A}}
\newcommand{\PP}{\mathbb{P}}
\newcommand{\E}{\mathbb{E}}
\newcommand{\eps}{\varepsilon}
\newcommand{\dd}{\,\mathrm{d}}
\begin{document}
\section*{UnsolvedMath ID: 10600019}
\textbf{Source problem code:} catalogue entry
\subsection{Definitions and mathematical statement}
For classical knots and an abelian coefficient group \(A\), a Vassiliev invariant extends to singular diagrams by
\[
 v(K_\times)=v(K_+)-v(K_-),
\]
and has finite type of degree at most \(d\) if this extension vanishes on diagrams with more than \(d\) singular double points. The source does not fix \(A\), orientation conventions, or the virtual finite-type formulation used in its key extension question.

\emph{Compound problem: Vassiliev invariants.} Understand the connection between virtual-knot polynomials and Vassiliev invariants of virtual knots in Kauffman's sense. The source reports prior work, but its citation and the family of polynomials to be related are absent.

The key question asks whether every classical finite-type invariant extends to a finite-type invariant of long virtual knots, specifically in the formulation of an unprovided reference. Let \(\mathcal K_{\mathrm{cl}}^{\mathrm{long}}\) be classical long knot classes, \(\mathcal K_{\mathrm{vir}}^{\mathrm{long}}\) long virtual knot classes, \(i\) the natural inclusion, and \(\operatorname{cl}\) classical closure. Once coefficients and the intended virtual filtration are fixed, its quantifiers are
\[
 \forall v\text{ of finite classical type},\quad
 \exists V:\mathcal K_{\mathrm{vir}}^{\mathrm{long}}\to A
 \text{ of finite virtual type},\quad
 V(i(L))=v(\operatorname{cl}(L))
 \quad\forall L\in\mathcal K_{\mathrm{cl}}^{\mathrm{long}}.
\]
The excerpt does not require the virtual extension to have the same degree as \(v\); if a degree-preserving statement is intended, it must be recovered from the cited formulation rather than added here.

The source also reports a natural alternative notion of Vassiliev invariants for virtual knots with a different finite-type condition, and asks that it be investigated. In common virtual-knot settings one can distinguish singularization using a positive-minus-negative classical crossing from a derivative using a classical-minus-virtual crossing. These can produce different filtrations, but the missing citations prevent unambiguously assigning the extension question to a particular one. The alternative is not identified with the first notion or silently substituted for it.
\subsection{Short English statement}
Relate virtual-knot polynomials to virtual Vassiliev invariants, and ask whether every classical finite-type invariant has some finite-type extension to long virtual knots. The source distinguishes competing virtual finite-type definitions; the missing reference is essential to choosing the intended one.
\subsection{Sources}
\begin{itemize}
\item Local UnsolvedMath registry, permanent ID 10600019, source code catalogue entry, ``Virtual-knot problem 19 --- Vassiliev Invariants.'' Kauffman is named for one sense of virtual Vassiliev invariants; no individual proposer is supplied.
\item \emph{Unsolved Problems in Virtual Knot Theory and Combinatorial Knot Theory}, revised survey, problem-list item 19. Supplied URL: \url{https://arxiv.org/abs/1409.2823}; recorded access 2026-07-29; extraction source-tex. The dedicated section contains 60 top-level enumerated items, with compound items retained intact.
\item Citations for prior polynomial connections, the exact extension formulation, and the alternative finite-type notion are missing from the extracted statement. No bibliographic details or identification of those references is invented.
\item Supplied status partially\_solved and PARTIAL-PROGRESS. The summary says the theory is well developed and special extension cases occur in the literature, but its author did not find a definitive full resolution; no specific theorem is supplied. Current primary-literature status and redistribution terms for the publicly accessible author/publisher copy remain NEEDS\_REVIEW.
\end{itemize}
\end{document}
